% Source: tex/chapters/ch06/03_導入：運用が重要な理由.tex
\subsection*{導入：運用が重要な理由}
\addcontentsline{toc}{subsection}{導入：運用が重要な理由}

ソフトウェアエンジニアリング運用とは、ソフトウェアまたはシステムの完全性と安定性を保ちながら、配備し、運用し、支援するために必要な活動と作業である。ここでいう運用は、対象環境へのソフトウェアの配備と設定、利用中の監視と管理、欠陥対応、環境変更への対応、新しい利用者要求への対応を含む。

従来、運用部門や計算センターは、開発部門から分離された組織として扱われることが多かった。この分離は、引き継ぎの遅れ、手作業による設定ミス、責任境界の曖昧化、リリースの遅延を生みやすい。現在は、開発、保守、運用を近づけ、DevOps、IaC、プラットフォーム工学、SREなどを使って、変更を早く、かつ安全に本番へ届ける考え方が広がっている。

この章で重要な視点は二つである。第一に、運用は「最後に引き受ける部門作業」ではなく、開発初期から設計されるサービスである。第二に、現代の運用では、環境、設定、配備、監視、試験、復旧をできるだけコードと自動化で表すことで、再現性、標準化、セキュリティ、変更管理、拡張性を高める。

\begin{pointbox}{重要}
運用の品質は、利用者が直接体験するソフトウェアの品質である。コードが正しくても、配備が失敗し、監視できず、復旧できず、容量が足りなければ、利用者から見たサービスは失敗である。
\end{pointbox}
